%==============================================================================
\section{Models --- Các loại mô hình và phương pháp học máy}
\label{sec:mo-hinh-ml}
%==============================================================================

\begin{dongco}[title={Nhiều cách xây mô hình từ dữ liệu}]
Có nhiều thuật toán, kỹ thuật và phương pháp ML để xây mô hình giải bài toán thực tế bằng dữ liệu.
\end{dongco}

\begin{vidu}[title={Tổng quan các kiểu học máy}]
\tsimg{11_mo_hinh_va_cac_kieu_hoc/img01.jpg}
\tsimg{11_mo_hinh_va_cac_kieu_hoc/img02.jpg}
\end{vidu}

\begin{phuongphap}[title={Bốn loại phương pháp ML (theo mức giám sát)}]
\begin{enumerate}
  \item Supervised Learning (học có giám sát)
  \item Unsupervised Learning (học không giám sát)
  \item Semi-supervised Learning (học bán giám sát)
  \item Reinforcement Learning (học tăng cường)
\end{enumerate}
\end{phuongphap}

\subsection{Supervised Learning (tổng quan)}

\begin{dinhnghia}[title={Học có giám sát}]
Thuật toán nhận tập huấn luyện gồm mẫu đầu vào và \textbf{nhãn/đáp án} (output) tương ứng.
target: học ánh xạ từ đầu vào $x$ sang đầu ra $Y$ sau nhiều lần huấn luyện:
\[
Y = f(x).
\]
Cần xấp xỉ $f$ tốt để dự đoán $Y$ cho dữ liệu $x$ mới.
\end{dinhnghia}

\begin{ynghia}[title={``Có giám sát''}]
Quá trình học như có ``giáo viên'' chỉ đáp án đúng.
Ví dụ thuật toán: Decision Tree, Random Forest, KNN, Logistic Regression, \ldots
\end{ynghia}

\begin{phuongphap}[title={Hai nhóm tác vụ supervised}]
\begin{itemize}
  \item \textbf{Classification} (phân loại)
  \item \textbf{Regression} (hồi quy)
\end{itemize}
\end{phuongphap}

\subsubsection{Classification}

\begin{ynghia}[title={target}]
Dự đoán nhãn phân loại (categorical) cho đầu vào; đầu ra là giá trị rời rạc, không có thứ tự bắt buộc, mỗi mẫu thuộc một lớp.
\end{ynghia}

\begin{phuongphap}[title={Một số mô hình phân loại phổ biến}]
Logistic Regression, Decision Trees, Random Forest, K-nearest Neighbor, Support Vector Machine, Naive Bayes, Linear Discriminant Analysis, Neural Networks.
\end{phuongphap}

\subsubsection{Regression}

\begin{ynghia}[title={target}]
Dự đoán giá trị số \textbf{liên tục}; mô hình học quan hệ giữa biến độc lập (feature) và biến phụ thuộc (outcome).
\end{ynghia}

\begin{phuongphap}[title={Một số mô hình hồi quy phổ biến}]
Linear Regression, Ridge Regression, Decision Trees, Random Forest, K-nearest Neighbor, Neural Network Regression.
\end{phuongphap}

\subsection{Unsupervised Learning (tổng quan)}

\begin{dinhnghia}[title={Học không giám sát}]
Ngược với supervised: \textbf{không} có nhãn/``giáo viên'' hướng dẫn.
Hữu ích khi không có dữ liệu đã gán nhãn sẵn nhưng cần trích pattern từ đầu vào $x$.
\end{dinhnghia}

\begin{ynghia}[title={Ví dụ trực giác}]
Chỉ có $x$, không có $Y$ tương ứng; thuật toán tự khám phá cấu trúc thú vị trong dữ liệu.
Ví dụ thuật toán: K-means clustering, K-nearest neighbors (trong ngữ cảnh unsupervised), \ldots
\end{ynghia}

\begin{phuongphap}[title={Ba nhóm tác vụ unsupervised}]
Clustering, Association, Dimensionality Reduction (và Anomaly Detection).
\end{phuongphap}

\subsubsection{Clustering}

\begin{ynghia}[title={Clustering}]
Tìm độ tương đồng/quan hệ giữa mẫu và gom thành nhóm theo feature.
Ví dụ thực tế: gom khách hàng theo hành vi mua hàng.
\end{ynghia}

\begin{phuongphap}[title={Mô hình clustering}]
K-Means, Hierarchical Clustering, Mean-shift, DBSCAN, HDBSCAN, BIRCH, Affinity Propagation, Agglomerative Clustering.
\end{phuongphap}

\subsubsection{Association}

\begin{ynghia}[title={Association / Market basket analysis}]
Phân tích tập dữ liệu lớn để tìm pattern thể hiện quan hệ giữa các mục (ví dụ sản phẩm hay mua cùng nhau).
\end{ynghia}

\begin{phuongphap}[title={Mô hình association}]
Apriori, Eclat, FP-growth.
\end{phuongphap}

\subsubsection{Dimensionality Reduction}

\begin{ynghia}[title={Giảm chiều}]
Chọn tập feature đại diện/chính để giảm số chiều mỗi mẫu --- tránh \textbf{curse of dimensionality} khi có hàng triệu feature.
\end{ynghia}

\begin{phuongphap}[title={Mô hình giảm chiều}]
PCA, Autoencoders, SVD.
\end{phuongphap}

\subsubsection{Anomaly Detection}

\begin{ynghia}[title={Phát hiện bất thường}]
Tìm sự kiện/quan sát hiếm, không thường xảy ra; phân biệt điểm bất thường và bình thường.
Một số thuật toán clustering, KNN có thể hỗ trợ phát hiện anomaly.
\end{ynghia}

\subsection{Semi-supervised Learning (tổng quan)}

\begin{dinhnghia}[title={Học bán giám sát}]
Không hoàn toàn supervised cũng không hoàn toàn unsupervised --- nằm giữa hai loại.
Thường dùng \textbf{ít} dữ liệu có nhãn và \textbf{nhiều} dữ liệu không nhãn.
\end{dinhnghia}

\begin{phuongphap}[title={Hai hướng triển khai}]
\begin{enumerate}
  \item Huấn luyện mô hình supervised trên ít dữ liệu có nhãn, áp dụng lên dữ liệu không nhãn để gán nhãn thêm, huấn luyện lại và lặp.
  \item Dùng unsupervised để cluster mẫu tương tự, gán nhãn cho cụm, rồi huấn luyện mô hình kết hợp thông tin này.
\end{enumerate}
\end{phuongphap}

\subsection{Reinforcement Learning (tổng quan)}

\begin{dinhnghia}[title={Học tăng cường}]
Khác các phương pháp trên và ít dùng hơn trong giới thiệu cơ bản.
Có \textbf{agent} được huấn luyện theo thời gian để tương tác môi trường: quan sát trạng thái, chọn hành động, nhận thưởng/phạt, cập nhật chiến lược.
\end{dinhnghia}

\begin{phuongphap}[title={Sáu bước}]
\begin{enumerate}
  \item Chuẩn bị agent với tập chiến lược ban đầu.
  \item Quan sát môi trường và trạng thái hiện tại.
  \item Chọn policy tối ưu theo trạng thái và thực hiện hành động.
  \item Nhận reward hoặc penalty.
  \item Cập nhật chiến lược nếu cần.
  \item Lặp bước 2--5 đến khi agent học được policy tối ưu.
\end{enumerate}
\end{phuongphap}

\begin{phuongphap}[title={Thuật toán RL phổ biến}]
Q-learning, Markov Decision Process (MDP), SARSA, DQN, DDPG.
\end{phuongphap}
